\begin{abstract}
Existing peptide--MHC models primarily predict binding or general presentation and therefore do not directly characterize tissue-dependent presentation preferences under a fixed MHC context. We formulate tissue--MHC-specific presentation preference as a multi-task prediction problem and construct paired human and mouse benchmarks in which positive and pseudo-negative peptides are matched by both MHC restriction and parent UniProt protein. We further develop the TissuePMHC human model, which combines a position-preserving multi-kernel peptide encoder with complementary global auxiliary and HLA-specific branches, fixed task-wise rank fusion, and multi-seed averaging. On the primary standard pair-disjoint benchmarks, this model achieves mean task AUROC/AUPRC values of 0.8448/0.8348 across 157 human tissue--HLA tasks, while a separately trained frozen five-seed Factorized MMoE achieves 0.8562/0.8506 across 24 mouse tissue--H2 tasks. Under connected-component peptide-disjoint out-of-fold evaluation, the frozen human and mouse models retain mean task AUROCs of 0.7652 and 0.7529, respectively, revealing a substantial unseen-peptide generalization gap. On the same strict folds, the human model is statistically comparable to the auxiliary dual-branch MLP (+0.0010 AUROC; \(q=0.442\)); the three mouse ensembles are likewise comparable. Capacity-matched controls that receive peptide and MHC but no tissue information reach strict AUROCs of 0.6630 and 0.6412, while frozen MHCflurry and NetMHCpan controls show moderate discrimination. Adding external scores by cross-fitted stacking improves AUROC by at most 0.0041. The results support strict task learnability and a retained human auxiliary-supervision benefit, not universal strict superiority of the full architectures. They establish a reproducible benchmark for tissue-conditioned MHC-I presentation ranking without implying unseen-task, protein-disjoint, external-cohort, or clinical generalization.
\end{abstract}
